\documentclass[10pt,a4paper]{amsart}
\usepackage{amsmath,amssymb,amsthm,graphicx}
% Retained preamble for the Stacks Project chapter snapshot in this directory.
% The chapter's own preamble.tex is a copy of the shared Stacks Project preamble; the block
% below is sliced MECHANICALLY, line for line, from it. Nothing was added, removed or
% redefined; the wrapper class, the xr-hyper/hyperref wiring and the externaldocument list
% are omitted because the assembler supplies the document class.
%
% lines taken from preamble.tex: 3, 5, 9, 11, 12, 156, 157, 158, ... (51 lines total)
%
% package coverage restored from the chapter's own preamble (2026-10-13): xy, xy-2cell

\usepackage[all]{xy}
\xyoption{2cell}
\UseAllTwocells

\newenvironment{reference}{\comment}{\endcomment}
\newenvironment{slogan}{\comment}{\endcomment}
\newenvironment{history}{\comment}{\endcomment}
\theoremstyle{plain}
\newtheorem{theorem}[subsection]{Theorem}
\newtheorem{proposition}[subsection]{Proposition}
\newtheorem{lemma}[subsection]{Lemma}
\theoremstyle{definition}
\newtheorem{definition}[subsection]{Definition}
\newtheorem{example}[subsection]{Example}
\newtheorem{exercise}[subsection]{Exercise}
\newtheorem{situation}[subsection]{Situation}
\theoremstyle{remark}
\newtheorem{remark}[subsection]{Remark}
\newtheorem{remarks}[subsection]{Remarks}
\numberwithin{equation}{subsection}
\def\lim{\mathop{\mathrm{lim}}\nolimits}
\def\colim{\mathop{\mathrm{colim}}\nolimits}
\def\Spec{\mathop{\mathrm{Spec}}}
\def\Hom{\mathop{\mathrm{Hom}}\nolimits}
\def\Ext{\mathop{\mathrm{Ext}}\nolimits}
\def\SheafHom{\mathop{\mathcal{H}\!\mathit{om}}\nolimits}
\def\SheafExt{\mathop{\mathcal{E}\!\mathit{xt}}\nolimits}
\def\Sch{\mathit{Sch}}
\def\Mor{\mathop{\mathrm{Mor}}\nolimits}
\def\Ob{\mathop{\mathrm{Ob}}\nolimits}
\def\Sh{\mathop{\mathit{Sh}}\nolimits}
\def\NL{\mathop{N\!L}\nolimits}
\def\CH{\mathop{\mathrm{CH}}\nolimits}
\def\proetale{{pro\text{-}\acute{e}tale}}
\def\etale{{\acute{e}tale}}
\def\QCoh{\mathit{QCoh}}
\def\Ker{\mathop{\mathrm{Ker}}}
\def\Im{\mathop{\mathrm{Im}}}
\def\Coker{\mathop{\mathrm{Coker}}}
\def\Coim{\mathop{\mathrm{Coim}}}
\DeclareMathSymbol{\boxtimes}{\mathbin}{AMSa}{"02}
\def\QCohstack{\mathcal{QC}\!\mathit{oh}}
\def\Cohstack{\mathcal{C}\!\mathit{oh}}
\def\Spacesstack{\mathcal{S}\!\mathit{paces}}
\def\Quotfunctor{\mathrm{Quot}}
\def\Hilbfunctor{\mathrm{Hilb}}
\def\Curvesstack{\mathcal{C}\!\mathit{urves}}
\def\Polarizedstack{\mathcal{P}\!\mathit{olarized}}
\def\Complexesstack{\mathcal{C}\!\mathit{omplexes}}
\def\Pic{\mathop{\mathrm{Pic}}\nolimits}
\def\Picardstack{\mathcal{P}\!\mathit{ic}}
\def\Picardfunctor{\mathrm{Pic}}
\def\Deformationcategory{\mathcal{D}\!\mathit{ef}}
\usepackage{cleveref}
\crefname{originallemma}{original lemma}{original lemmas}
\Crefname{originallemma}{Original lemma}{Original lemmas}
\makeatletter
\newlabel{algebra-lemma-integral-overring-surjective}{{00GQ}{source}{}{sourcepointer-1}{}}
\newlabel{algebra-lemma-integral-overring-surjective@cref}{{[originallemma][0][]00GQ}{[1][0][]source}{}{}{}}
\makeatother
\crefname{originallemma}{original lemma}{original lemmas}
\Crefname{originallemma}{Original lemma}{Original lemmas}
\makeatletter
\newlabel{lemma-colimit-check-etale-locally}{{0GHL}{source}{}{sourcepointer-2}{}}
\newlabel{lemma-colimit-check-etale-locally@cref}{{[originallemma][0][]0GHL}{[1][0][]source}{}{}{}}
\makeatother
\crefname{originallemma}{original lemma}{original lemmas}
\Crefname{originallemma}{Original lemma}{Original lemmas}
\makeatletter
\newlabel{lemma-descend-etale-proper-surjective}{{0GG0}{source}{}{sourcepointer-3}{}}
\newlabel{lemma-descend-etale-proper-surjective@cref}{{[originallemma][0][]0GG0}{[1][0][]source}{}{}{}}
\makeatother
\crefname{originallemma}{original lemma}{original lemmas}
\Crefname{originallemma}{Original lemma}{Original lemmas}
\makeatletter
\newlabel{spaces-cohomology-lemma-flat-base-change-cohomology}{{073K}{source}{}{sourcepointer-4}{}}
\newlabel{spaces-cohomology-lemma-flat-base-change-cohomology@cref}{{[originallemma][0][]073K}{[1][0][]source}{}{}{}}
\makeatother
\crefname{originallemma}{original lemma}{original lemmas}
\Crefname{originallemma}{Original lemma}{Original lemmas}
\makeatletter
\newlabel{spaces-cohomology-lemma-proper-pushforward-coherent}{{08AR}{source}{}{sourcepointer-5}{}}
\newlabel{spaces-cohomology-lemma-proper-pushforward-coherent@cref}{{[originallemma][0][]08AR}{[1][0][]source}{}{}{}}
\makeatother
\crefname{originallemma}{original lemma}{original lemmas}
\Crefname{originallemma}{Original lemma}{Original lemmas}
\makeatletter
\newlabel{spaces-more-morphisms-theorem-stein-factorization-Noetherian}{{0A1B}{source}{}{sourcepointer-6}{}}
\newlabel{spaces-more-morphisms-theorem-stein-factorization-Noetherian@cref}{{[originallemma][0][]0A1B}{[1][0][]source}{}{}{}}
\makeatother
\crefname{originallemma}{original lemma}{original lemmas}
\Crefname{originallemma}{Original lemma}{Original lemmas}
\makeatletter
\newlabel{spaces-morphisms-definition-relative-spec}{{081W}{source}{}{sourcepointer-7}{}}
\newlabel{spaces-morphisms-definition-relative-spec@cref}{{[originallemma][0][]081W}{[1][0][]source}{}{}{}}
\makeatother
\crefname{originallemma}{original lemma}{original lemmas}
\Crefname{originallemma}{Original lemma}{Original lemmas}
\makeatletter
\newlabel{spaces-morphisms-definition-universally-injective}{{03MV}{source}{}{sourcepointer-8}{}}
\newlabel{spaces-morphisms-definition-universally-injective@cref}{{[originallemma][0][]03MV}{[1][0][]source}{}{}{}}
\makeatother
\crefname{originallemma}{original lemma}{original lemmas}
\Crefname{originallemma}{Original lemma}{Original lemmas}
\makeatletter
\newlabel{spaces-morphisms-lemma-base-change-proper}{{04WP}{source}{}{sourcepointer-9}{}}
\newlabel{spaces-morphisms-lemma-base-change-proper@cref}{{[originallemma][0][]04WP}{[1][0][]source}{}{}{}}
\makeatother
\crefname{originallemma}{original lemma}{original lemmas}
\Crefname{originallemma}{Original lemma}{Original lemmas}
\makeatletter
\newlabel{spaces-morphisms-lemma-base-change-surjective}{{03MH}{source}{}{sourcepointer-10}{}}
\newlabel{spaces-morphisms-lemma-base-change-surjective@cref}{{[originallemma][0][]03MH}{[1][0][]source}{}{}{}}
\makeatother
\crefname{originallemma}{original lemma}{original lemmas}
\Crefname{originallemma}{Original lemma}{Original lemmas}
\makeatletter
\newlabel{spaces-morphisms-lemma-composition-proper}{{04XY}{source}{}{sourcepointer-11}{}}
\newlabel{spaces-morphisms-lemma-composition-proper@cref}{{[originallemma][0][]04XY}{[1][0][]source}{}{}{}}
\makeatother
\crefname{originallemma}{original lemma}{original lemmas}
\Crefname{originallemma}{Original lemma}{Original lemmas}
\makeatletter
\newlabel{spaces-morphisms-lemma-composition-surjective}{{03MG}{source}{}{sourcepointer-12}{}}
\newlabel{spaces-morphisms-lemma-composition-surjective@cref}{{[originallemma][0][]03MG}{[1][0][]source}{}{}{}}
\makeatother
\crefname{originallemma}{original lemma}{original lemmas}
\Crefname{originallemma}{Original lemma}{Original lemmas}
\makeatletter
\newlabel{spaces-morphisms-lemma-locally-finite-type-locally-noetherian}{{04ZK}{source}{}{sourcepointer-13}{}}
\newlabel{spaces-morphisms-lemma-locally-finite-type-locally-noetherian@cref}{{[originallemma][0][]04ZK}{[1][0][]source}{}{}{}}
\makeatother
\crefname{originallemma}{original lemma}{original lemmas}
\Crefname{originallemma}{Original lemma}{Original lemmas}
\makeatletter
\newlabel{spaces-morphisms-lemma-universally-closed-permanence}{{04NX}{source}{}{sourcepointer-14}{}}
\newlabel{spaces-morphisms-lemma-universally-closed-permanence@cref}{{[originallemma][0][]04NX}{[1][0][]source}{}{}{}}
\makeatother
\crefname{originallemma}{original lemma}{original lemmas}
\Crefname{originallemma}{Original lemma}{Original lemmas}
\makeatletter
\newlabel{spaces-morphisms-lemma-universally-injective}{{040X}{source}{}{sourcepointer-15}{}}
\newlabel{spaces-morphisms-lemma-universally-injective@cref}{{[originallemma][0][]040X}{[1][0][]source}{}{}{}}
\makeatother
\crefname{originallemma}{original lemma}{original lemmas}
\Crefname{originallemma}{Original lemma}{Original lemmas}
\makeatletter
\newlabel{spaces-morphisms-remark-factorization-quasi-compact-quasi-separated}{{081X}{source}{}{sourcepointer-16}{}}
\newlabel{spaces-morphisms-remark-factorization-quasi-compact-quasi-separated@cref}{{[originallemma][0][]081X}{[1][0][]source}{}{}{}}
\makeatother
\crefname{originallemma}{original lemma}{original lemmas}
\Crefname{originallemma}{Original lemma}{Original lemmas}
\makeatletter
\newlabel{spaces-properties-lemma-morphism-to-affine-scheme}{{05Z1}{source}{}{sourcepointer-17}{}}
\newlabel{spaces-properties-lemma-morphism-to-affine-scheme@cref}{{[originallemma][0][]05Z1}{[1][0][]source}{}{}{}}
\makeatother
\title{Coequalizer of a proper groupoid on a locally Noetherian algebraic space exists and is finite}
\author{The Stacks Project Authors}
\date{stacks-project:spaces-pushouts.tex:local-snapshot-20261012; GNU FDL 1.2 or later}
\begin{document}
\maketitle
\noindent\textit{Editorial note.} Statement span is lines 2862-2880 of sources/stacks-spaces-pushouts/pinned-original-source.tex; proof span is lines 2882-3015, the complete author proof (132 body lines).\par\medskip
\noindent\textit{Editorial note.} Navigation stubs: 17 cited-but-not-included labels are supplied as non-load-bearing pointers with their pinned tags where the snapshot tag index has them.\par\medskip
\noindent\textit{Editorial note.} extra\_packages is empty on purpose: the retained preamble already loads xy with options and a second declaration clashes in the assembler wrapper.\par\medskip
\noindent\textit{Editorial note.} Text is retained verbatim, including the author's own typos and placeholders.\par\medskip
\noindent\textit{Editorial note.} Queue index 1 of queue\_spaces-pushouts.json (pb=5812, stmt 2862-2880, proof 2882-3015).\par\medskip
\noindent\textit{Editorial note.} spaces-pushouts wave-1: fresh adjudication accept (2026-10-14); six-check dispatched pool (§2 row appended 02:4x).\par\medskip
\noindent\textit{Editorial note.} Cross-chapter same-label warning (NOT an occupancy conflict, Chief ruling 2026-10-13 evening): label 'lemma-coequalizer' also exists in stacks-snapshot-20261011 — a different proposition in a different chapter file.\par\medskip
\noindent\textit{Allowed standard prerequisites (editorial).}\begin{itemize}
\item Explicit prerequisite boundary. The cited inputs are prerequisites; the argument in the retained span is not replaced by any cited result.
\end{itemize}\medskip
\ifdefined\hypertarget\hypertarget{sourcepointer-1}{}\fi
\noindent\textit{Source reference.} Original lemma 00GQ: sources/stacks-spaces-pushouts/pinned-original-source.tex, source lines 3015--3015. Cited by the retained statement or proof but not part of the collected spans; supplied only as a navigation pointer. The argument does not depend on its content.\par\medskip
\ifdefined\hypertarget\hypertarget{sourcepointer-2}{}\fi
\noindent\textit{Source reference.} Original lemma 0GHL: sources/stacks-spaces-pushouts/pinned-original-source.tex, source lines 450--469. Cited by the retained statement or proof but not part of the collected spans; supplied only as a navigation pointer. The argument does not depend on its content.\par\medskip
\ifdefined\hypertarget\hypertarget{sourcepointer-3}{}\fi
\noindent\textit{Source reference.} Original lemma 0GG0: sources/stacks-spaces-pushouts/pinned-original-source.tex, source lines 966--974. Cited by the retained statement or proof but not part of the collected spans; supplied only as a navigation pointer. The argument does not depend on its content.\par\medskip
\ifdefined\hypertarget\hypertarget{sourcepointer-4}{}\fi
\noindent\textit{Source reference.} Original lemma 073K: sources/stacks-spaces-pushouts/pinned-original-source.tex, source lines 3015--3015. Cited by the retained statement or proof but not part of the collected spans; supplied only as a navigation pointer. The argument does not depend on its content.\par\medskip
\ifdefined\hypertarget\hypertarget{sourcepointer-5}{}\fi
\noindent\textit{Source reference.} Original lemma 08AR: sources/stacks-spaces-pushouts/pinned-original-source.tex, source lines 3015--3015. Cited by the retained statement or proof but not part of the collected spans; supplied only as a navigation pointer. The argument does not depend on its content.\par\medskip
\ifdefined\hypertarget\hypertarget{sourcepointer-6}{}\fi
\noindent\textit{Source reference.} Original lemma 0A1B: sources/stacks-spaces-pushouts/pinned-original-source.tex, source lines 3015--3015. Cited by the retained statement or proof but not part of the collected spans; supplied only as a navigation pointer. The argument does not depend on its content.\par\medskip
\ifdefined\hypertarget\hypertarget{sourcepointer-7}{}\fi
\noindent\textit{Source reference.} Original lemma 081W: sources/stacks-spaces-pushouts/pinned-original-source.tex, source lines 3015--3015. Cited by the retained statement or proof but not part of the collected spans; supplied only as a navigation pointer. The argument does not depend on its content.\par\medskip
\ifdefined\hypertarget\hypertarget{sourcepointer-8}{}\fi
\noindent\textit{Source reference.} Original lemma 03MV: sources/stacks-spaces-pushouts/pinned-original-source.tex, source lines 3015--3015. Cited by the retained statement or proof but not part of the collected spans; supplied only as a navigation pointer. The argument does not depend on its content.\par\medskip
\ifdefined\hypertarget\hypertarget{sourcepointer-9}{}\fi
\noindent\textit{Source reference.} Original lemma 04WP: sources/stacks-spaces-pushouts/pinned-original-source.tex, source lines 3015--3015. Cited by the retained statement or proof but not part of the collected spans; supplied only as a navigation pointer. The argument does not depend on its content.\par\medskip
\ifdefined\hypertarget\hypertarget{sourcepointer-10}{}\fi
\noindent\textit{Source reference.} Original lemma 03MH: sources/stacks-spaces-pushouts/pinned-original-source.tex, source lines 3015--3015. Cited by the retained statement or proof but not part of the collected spans; supplied only as a navigation pointer. The argument does not depend on its content.\par\medskip
\ifdefined\hypertarget\hypertarget{sourcepointer-11}{}\fi
\noindent\textit{Source reference.} Original lemma 04XY: sources/stacks-spaces-pushouts/pinned-original-source.tex, source lines 3015--3015. Cited by the retained statement or proof but not part of the collected spans; supplied only as a navigation pointer. The argument does not depend on its content.\par\medskip
\ifdefined\hypertarget\hypertarget{sourcepointer-12}{}\fi
\noindent\textit{Source reference.} Original lemma 03MG: sources/stacks-spaces-pushouts/pinned-original-source.tex, source lines 3015--3015. Cited by the retained statement or proof but not part of the collected spans; supplied only as a navigation pointer. The argument does not depend on its content.\par\medskip
\ifdefined\hypertarget\hypertarget{sourcepointer-13}{}\fi
\noindent\textit{Source reference.} Original lemma 04ZK: sources/stacks-spaces-pushouts/pinned-original-source.tex, source lines 3015--3015. Cited by the retained statement or proof but not part of the collected spans; supplied only as a navigation pointer. The argument does not depend on its content.\par\medskip
\ifdefined\hypertarget\hypertarget{sourcepointer-14}{}\fi
\noindent\textit{Source reference.} Original lemma 04NX: sources/stacks-spaces-pushouts/pinned-original-source.tex, source lines 3015--3015. Cited by the retained statement or proof but not part of the collected spans; supplied only as a navigation pointer. The argument does not depend on its content.\par\medskip
\ifdefined\hypertarget\hypertarget{sourcepointer-15}{}\fi
\noindent\textit{Source reference.} Original lemma 040X: sources/stacks-spaces-pushouts/pinned-original-source.tex, source lines 3015--3015. Cited by the retained statement or proof but not part of the collected spans; supplied only as a navigation pointer. The argument does not depend on its content.\par\medskip
\ifdefined\hypertarget\hypertarget{sourcepointer-16}{}\fi
\noindent\textit{Source reference.} Original lemma 081X: sources/stacks-spaces-pushouts/pinned-original-source.tex, source lines 3015--3015. Cited by the retained statement or proof but not part of the collected spans; supplied only as a navigation pointer. The argument does not depend on its content.\par\medskip
\ifdefined\hypertarget\hypertarget{sourcepointer-17}{}\fi
\noindent\textit{Source reference.} Original lemma 05Z1: sources/stacks-spaces-pushouts/pinned-original-source.tex, source lines 3015--3015. Cited by the retained statement or proof but not part of the collected spans; supplied only as a navigation pointer. The argument does not depend on its content.\par\medskip
\begin{lemma}
\label{lemma-coequalizer}
Let $S$ be a scheme. Let
$$
g : Y \longrightarrow X
$$
be a morphism of algebraic spaces over $S$. Assume $X$ is locally Noetherian,
and $g$ is proper. Let $R = Y \times_X Y$ with projection morphisms
$t, s : R \to Y$. There exists a coequalizer $X'$ of $s, t : R \to Y$
in the category of algebraic spaces over $S$. Moreover
\begin{enumerate}
\item The morphism $X' \to X$ is finite.
\item The morphism $Y \to X'$ is proper.
\item The morphism $Y \to X'$ is surjective.
\item The morphism $X' \to X$ is universally injective.
\item If $g$ is surjective, the morphism $X' \to X$
is a universal homeomorphism.
\end{enumerate}
\end{lemma}

\begin{proof}
Denote $h : R \to X$ denote the composition of either $s$ or $t$
with $g$. Then $h$ is proper by Morphisms of Spaces, Lemmas
\ref{spaces-morphisms-lemma-base-change-proper} and
\ref{spaces-morphisms-lemma-composition-proper}.
The sheaves
$$
g_*\mathcal{O}_Y
\quad\text{and}\quad
h_*\mathcal{O}_R
$$
are coherent $\mathcal{O}_X$-algebras by Cohomology of Spaces, Lemma
\ref{spaces-cohomology-lemma-proper-pushforward-coherent}.
The $X$-morphisms $s$, $t$ induce $\mathcal{O}_X$-algebra maps
$s^\sharp, t^\sharp$ from the first to the second.
Set
$$
\mathcal{A} = \text{Equalizer}\left(s^\sharp, t^\sharp :
g_*\mathcal{O}_Y \longrightarrow h_*\mathcal{O}_R\right)
$$
Then $\mathcal{A}$ is a coherent $\mathcal{O}_X$-algebra and we
can define
$$
X' = \underline{\Spec}_X(\mathcal{A})
$$
as in Morphisms of Spaces, Definition
\ref{spaces-morphisms-definition-relative-spec}.
By Morphisms of Spaces, Remark
\ref{spaces-morphisms-remark-factorization-quasi-compact-quasi-separated}
and functoriality of the $\underline{\Spec}$ construction
there is a factorization
$$
Y \longrightarrow X' \longrightarrow X
$$
and the morphism $g' : Y \to X'$ equalizes $s$ and $t$.

\medskip\noindent
Before we show that $X'$ is the coequalizer of $s$ and $t$, we show
that $Y \to X'$ and $X' \to X$ have the desired properties. Since $\mathcal{A}$
is a coherent $\mathcal{O}_X$-module it is clear that $X' \to X$ is a
finite morphism of algebraic spaces. This proves (1).
The morphism $Y \to X'$ is proper by Morphisms of Spaces, Lemma
\ref{spaces-morphisms-lemma-universally-closed-permanence}. This proves (2).
Denote $Y \to Y' \to X$ with
$Y' = \underline{\Spec}_X(g_*\mathcal{O}_Y)$ the Stein factorization of $g$,
see More on Morphisms of Spaces, Theorem
\ref{spaces-more-morphisms-theorem-stein-factorization-Noetherian}.
Of course we obtain morphisms $Y \to Y' \to X' \to X$ fitting
with the morphisms studied above.
Since $\mathcal{O}_{X'} \subset g_*\mathcal{O}_Y$ is a finite extension
we see that $Y' \to X'$ is finite and surjective.
Some details omitted; hint: use
Algebra, Lemma \ref{algebra-lemma-integral-overring-surjective}
and reduce to the affine case by \'etale localization.
Since $Y \to Y'$ is surjective (with geometrically connected fibres)
we conclude that $Y \to X'$ is surjective. This proves (3).
To show that $X' \to X$ is universally injective, we have to show
that $X' \to X' \times_X X'$ is surjective, see
Morphisms of Spaces, Definition
\ref{spaces-morphisms-definition-universally-injective} and
Lemma \ref{spaces-morphisms-lemma-universally-injective}.
Since $Y \to X'$ is surjective (see above) and since base changes
and compositions of surjective morphisms are surjective by
Morphisms of Spaces, Lemmas
\ref{spaces-morphisms-lemma-base-change-surjective} and
\ref{spaces-morphisms-lemma-composition-surjective}
we see that $Y \times_X Y \to X' \times_X X'$ is surjective.
However, since $Y \to X'$ equalizes $s$ and $t$, we see that
$Y \times_X Y \to X' \times_X X'$ factors through $X' \to X' \times_X X'$
and we conclude this latter map is surjective. This proves (4).
Finally, if $g$ is surjective, then since $g$ factors through $X' \to X$
we see that $X' \to X$ is surjective. Since a surjective, universally
injective, finite morphism is a universal homeomorphism (because it
is universally bijective and universally closed), this proves (5).

\medskip\noindent
In the rest of the proof we show that $Y \to X'$ is the coequalizer
of $s$ and $t$ in the category of algebraic spaces over $S$.
Observe that $X'$ is locally Noetherian
(Morphisms of Spaces, Lemma
\ref{spaces-morphisms-lemma-locally-finite-type-locally-noetherian}).
Moreover, observe that $Y \times_{X'} Y \to Y \times_X Y$ is an
isomorphism as $Y \to X'$ equalizes $s$ and $t$ (this is a categorical
statement). Hence in order to prove the statement that $Y \to X'$
is the coequalizer of $s$ and $t$, we may and do assume $X = X'$.
In other words, $\mathcal{O}_X$ is the equalizer of the maps
$s^\sharp, t^\sharp : g_*\mathcal{O}_Y \to h_*\mathcal{O}_R$.

\medskip\noindent
Let $X_1 \to X$ be a flat morphism of algebraic spaces over $S$ with $X_1$
locally Noetherian. Denote $g_1 : Y_1 \to X_1$, $h_1 : R_1 \to X_1$ and
$s_1, t_1 : R_1 \to Y_1$ the base changes of $g, h, s, t$ to $X_1$.
Of course $g_1$ is proper and $R_1 = Y_1 \times_{X_1} Y_1$.
Since we have flat base change for pushforward of quasi-coherent modules,
Cohomology of Spaces, Lemma
\ref{spaces-cohomology-lemma-flat-base-change-cohomology}, we see that
$\mathcal{O}_{X_1}$ is the equalizer of the maps
$s_1^\sharp, t_1^\sharp : g_{1, *}\mathcal{O}_{Y_1} \to
h_{1, *}\mathcal{O}_{R_1}$. Hence all the assumptions we have
are preserved by this base change.

\medskip\noindent
At this point we are going to check conditions (1) and (2) of
Lemma \ref{lemma-colimit-check-etale-locally}. Condition (1)
follows from Lemma \ref{lemma-descend-etale-proper-surjective}
and the fact that $g$ is proper and surjective (because $X = X'$).
To check condition (2), by the remarks on base change above,
we reduce to the statement discussed and proved in the next paragraph.

\medskip\noindent
Assume $S = \Spec(A)$ is an affine scheme, $X = X'$ is an affine scheme, and
$Z$ is an affine scheme over $S$. We have to show that
$$
\Mor_S(X, Z) \longrightarrow
\text{Equalizer}(s, t : \Mor_S(Y, Z) \to \Mor_S(R, Z))
$$
is bijective. However, this is clear from the fact that $X = X'$
which implies $\mathcal{O}_X$ is the equalizer of the maps
$s^\sharp, t^\sharp : g_*\mathcal{O}_Y \to h_*\mathcal{O}_R$
which in turn implies
$$
\Gamma(X, \mathcal{O}_X) =
\text{Equalizer}\left(
s^\sharp, t^\sharp : \Gamma(Y, \mathcal{O}_Y) \to
\Gamma(R, \mathcal{O}_R)
\right)
$$
Namely, we have
$$
\Mor_S(X, Z) = \Hom_A(\Gamma(Z, \mathcal{O}_Z), \Gamma(X, \mathcal{O}_X))
$$
and similarly for $Y$ and $R$, see Properties of Spaces, Lemma
\ref{spaces-properties-lemma-morphism-to-affine-scheme}.
\end{proof}

\end{document}
